\section{Discovery changes the action space}\label{sec:C18:discovery}
Discovery provides the agent or harness with the currently available capabilities. It can involve a static catalog, server-provided schemas, search over tool descriptions, or staged loading of detailed definitions. The catalog is part of the policy's information state. Adding or removing tools changes the represented proposal interface and can alter performance even when the model weights and task distribution remain fixed. Discovery does not itself grant authority: the dispatcher independently restricts which discovered operations may execute.

Large catalogs create both token and discrimination costs. Injecting every complete schema can consume context and expose near-duplicate actions. A two-stage mechanism can retrieve candidate tools using coarse descriptions and load exact schemas only before invocation. This reduces active input but can exclude the necessary tool at the discovery stage. Evaluation should report catalog recall separately from correct use among discovered candidates. An oracle candidate list is a diagnostic condition, not an executable baseline.

Tool descriptions are themselves data that can be stale, misleading, or adversarial. A server can expose a harmless-looking title while its implementation has broader effects. A retrieved description can include instructions unrelated to the capability. Treating the description as an authoritative system instruction expands the trust boundary. Metadata can guide selection without being allowed to redefine the user's task or authority.

Current MCP separates discovery and invocation through its tool feature and schemas~\cite{mcp2026tools}. A capability catalog can vary with per-request authorization, so a cached list should be associated with the authorization context that produced it. The detailed protocol semantics are developed in Chapter~\ref{ch:C19}. The application-level consequence is that catalog identity, credential scope, and schema version belong in the decision record. A tool visible to one principal need not be visible to another.

Discovery can also create a training-deployment mismatch. A policy trained with a fixed tool ordering can exploit names or positions rather than robust semantic selection. At deployment, reordered or expanded catalogs can degrade selection. Controlled tests should permute catalog order, introduce unrelated near-name tools, and hold schema content constant. A model that fails these tests may have learned a brittle interface shortcut rather than a transferable tool-selection mechanism.

The proposed implementation caches catalogs with explicit expiration and invalidation, stores a schema digest per selected tool, and validates the exact schema again before dispatch where required. It distinguishes a missing capability from an unavailable credential and from a temporary execution failure. Those errors lead to different next actions: revise the plan, request an authorized scope change, or retry under a bounded recovery policy.

